\section{Результаты экспериментов с программной системой}
\subsection{Анализ модуля распределения заказов}

В разработанной интеллектуальной системе реализованы четыре метода построения расписаний: быстрая эвристика EDD+LWKR, нейросетевой диспетчер PointerNet (Behavioral Cloning), модель SLIM (Self‑Labeling) и PPO-агент (обучение с подкреплением). Для объективного сравнения был сформирован тестовый сценарий из 120 заказов трёх типов (фрезеровка, сборка, сварка) с пятью ресурсами, различающимися готовностью (1.0, 0.9, 0.7, 0.8, 0.3) и ограничениями по типам операций. Сроки сдачи заказов задавались жёсткими (3–12 часов), а длительности операций составляли 40–100 минут. Все эксперименты проводились при фиксированном начальном состоянии генератора псевдослучайных чисел (seed = 42), что обеспечивает полную воспроизводимость результатов.

\subsubsection{Сравнение среднего опоздания заказов}

На рисунке ~\ref{avg_types:image} приведена диаграмма среднего опоздания (в часах) для каждого метода. Опоздание определялось как разница между временем завершения операции и сроком сдачи заказа, усреднённая по всем 120 заказам. Чем меньше значение, тем лучше метод справляется с соблюдением сроков.

\begin{figure}[ht]
\includegraphics[width=1\linewidth]{avg_types}
\caption{Среднее опоздание заказов по методам планирования}
\label{avg_types:image}
\end{figure}

Из полученных данных видно, что наименьшее среднее опоздание достигается при использовании модели SLIM (8.1 часа), далее следуют PointerNet и эвристика (по 9.0 часов). Наибольшее опоздание зафиксировано у PPO-агента (13.3 часа), что объясняется использованием случайных признаков в процессе обучения. Таким образом, нейросетевые методы, обученные на качественных синтетических данных, способны не только не уступать эвристике, но и превосходить её по критерию соблюдения сроков.

\subsubsection{Анализ длительности расписания (makespan)}
